\section*{Rangkuman}
\addcontentsline{toc}{section}{Rangkuman}

Pada bab ini telah dibahas:
\begin{itemize}[leftmargin=*]
  \item keluarga translator dan perbedaan compiler, interpreter,
        assembler, serta cross-compiler;
  \item fase analisis (leksikal, sintaksis, semantik) dan sintesis
        (IR, optimasi, pembangkitan kode);
  \item peran tabel simbol dan strategi penanganan kesalahan;
  \item penelusuran fase: token, pohon sintaks, dan kode antara;
  \item cara membaca dan memeriksa keluaran antara setiap fase.
\end{itemize}

Bab IV memperdalam landasan formal yang dipakai pada setiap fase,
khususnya hierarki Chomsky dan hubungan antara kelas bahasa dengan
mesin pengenal. Pemahaman ini akan menjelaskan mengapa lexer cukup
memakai otomata hingga, sedangkan parser memerlukan otomata pushdown.
